\section{Java API}

\begin{question}
	以下哪个方法可以将字符串转换为整数？
	\begin{options}
		\item parseInt()
		\item valueOf()
		\item toString()
		\item toInteger()
	\end{options}
	\begin{answer}
		A
	\end{answer}
	\begin{explanation}
		\texttt{Integer.parseInt(String s)} 将字符串解析为 \texttt{int} 类型。
	\end{explanation}
\end{question}

\begin{question}
	在Java中，以下哪个方法用于获取字符串的长度？
	\begin{options}
		\item length()
		\item size()
		\item count()
		\item getLength()
	\end{options}
	\begin{answer}
		A
	\end{answer}
	\begin{explanation}
		字符串使用 \texttt{length()} 方法获取长度。
	\end{explanation}
\end{question}

\begin{question}
	以下哪个方法用于将字符串转换为大写？
	\begin{options}
		\item toUpperCase()
		\item upperCase()
		\item makeUpperCase()
		\item convertToUpperCase()
	\end{options}
	\begin{answer}
		A
	\end{answer}
	\begin{explanation}
		\texttt{toUpperCase()} 将字符串转为大写。
	\end{explanation}
\end{question}

\begin{question}
	以下哪个方法用于将字符串转换为小写？
	\begin{options}
		\item toLowerCase()
		\item lowerCase()
		\item makeLowerCase()
		\item convertToLowerCase()
	\end{options}
	\begin{answer}
		A
	\end{answer}
	\begin{explanation}
		\texttt{toLowerCase()} 将字符串转为小写。
	\end{explanation}
\end{question}

\begin{question}
	以下哪个方法用于截取字符串？
	\begin{options}
		\item substring()
		\item cut()
		\item slice()
		\item truncate()
	\end{options}
	\begin{answer}
		A
	\end{answer}
	\begin{explanation}
		\texttt{substring(int beginIndex)} 或 \texttt{substring(int begin, int end)} 用于截取。
	\end{explanation}
\end{question}

\begin{question}
	以下哪个方法用于判断字符串是否以指定字符串开头？
	\begin{options}
		\item startsWith()
		\item beginWith()
		\item isStartWith()
		\item hasStartWith()
	\end{options}
	\begin{answer}
		A
	\end{answer}
	\begin{explanation}
		\texttt{startsWith(String prefix)} 判断是否以某字符串开头。
	\end{explanation}
\end{question}

\begin{question}
	以下哪个方法用于判断字符串是否以指定字符串结尾？
	\begin{options}
		\item endsWith()
		\item finishWith()
		\item isEndWith()
		\item hasEndWith()
	\end{options}
	\begin{answer}
		A
	\end{answer}
	\begin{explanation}
		\texttt{endsWith(String suffix)} 判断是否以某字符串结尾。
	\end{explanation}
\end{question}

\begin{question}
	以下哪个方法用于在字符串中查找指定字符第一次出现的位置？
	\begin{options}
		\item indexOf()
		\item find()
		\item search()
		\item locate()
	\end{options}
	\begin{answer}
		A
	\end{answer}
	\begin{explanation}
		\texttt{indexOf()} 返回首次出现的位置，未找到返回 -1。
	\end{explanation}
\end{question}

\begin{question}
	以下哪个方法用于在字符串中查找指定字符最后一次出现的位置？
	\begin{options}
		\item lastIndexOf()
		\item findLast()
		\item searchLast()
		\item locateLast()
	\end{options}
	\begin{answer}
		A
	\end{answer}
	\begin{explanation}
		\texttt{lastIndexOf()} 返回最后一次出现的位置。
	\end{explanation}
\end{question}

\begin{question}
	以下关于Java中的包装类的说法错误的是
	\begin{options}
		\item 包装类可以将基本数据类型转换为对象
		\item Integer是int的包装类
		\item 包装类的对象不可变
		\item 包装类可以提高基本数据类型的操作效率
	\end{options}
	\begin{answer}
		D
	\end{answer}
	\begin{explanation}
		包装类主要用于对象化基本类型，便于集合操作，但不提高基本类型的操作效率。
	\end{explanation}
\end{question}

\begin{question}
	以下哪个是Double类的静态方法，用于将字符串转换为Double类型？
	\begin{options}
		\item parseDouble()
		\item valueOf()
		\item toString()
		\item toDouble()
	\end{options}
	\begin{answer}
		A
	\end{answer}
	\begin{explanation}
		\texttt{Double.parseDouble(String s)} 返回 \texttt{double} 值。
	\end{explanation}
\end{question}

\begin{question}
	以下关于Java中的枚举类型的说法错误的是
	\begin{options}
		\item 枚举类型是一种特殊的类
		\item 枚举类型的成员是常量
		\item 枚举类型可以有构造方法
		\item 枚举类型不能实现接口
	\end{options}
	\begin{answer}
		D
	\end{answer}
	\begin{explanation}
		枚举可以实现接口。
	\end{explanation}
\end{question}

\begin{question}
	以下哪个类用于获取系统属性？
	\begin{options}
		\item System
		\item Properties
		\item Runtime
		\item SystemProperties
	\end{options}
	\begin{answer}
		B
	\end{answer}
	\begin{explanation}
		\texttt{System.getProperties()} 返回 \texttt{Properties} 对象，用于获取系统属性。
	\end{explanation}
\end{question}

\begin{question}
	以下关于Java中的注解的说法错误的是
	\begin{options}
		\item 注解是一种元数据
		\item 可以自定义注解
		\item 注解可以影响程序的逻辑
		\item 注解可以为代码添加额外的信息
	\end{options}
	\begin{answer}
		C
	\end{answer}
	\begin{explanation}
		注解本身不改变程序逻辑，需由工具或框架处理。
	\end{explanation}
\end{question}

\begin{question}
	以下哪个是Java中的基本注解？
	\begin{options}
		\item @Override
		\item @Deprecated
		\item @SuppressWarnings
		\item 以上都是
	\end{options}
	\begin{answer}
		D
	\end{answer}
	\begin{explanation}
		这三个都是Java内置的基本注解。
	\end{explanation}
\end{question}

\begin{question}
	以下关于Java注解@SuppressWarnings的作用，正确的是
	\begin{options}
		\item 抑制所有的警告
		\item 抑制指定类型的警告
		\item 增加警告信息
		\item 用于代码注释
	\end{options}
	\begin{answer}
		B
	\end{answer}
	\begin{explanation}
		可以指定警告类型，如 \texttt{"unchecked"}。
	\end{explanation}
\end{question}

\begin{question}
	以下关于Java中StringBuilder和StringBuffer的说法正确的是
	\begin{options}
		\item StringBuilder是线程安全的，StringBuffer不是
		\item StringBuffer是线程安全的，StringBuilder不是
		\item 两者性能相同
		\item 两者都不是线程安全的
	\end{options}
	\begin{answer}
		B
	\end{answer}
	\begin{explanation}
		\texttt{StringBuffer} 方法使用 \texttt{synchronized} 修饰，线程安全。
	\end{explanation}
\end{question}

\begin{question}
	以下哪个方法用于将StringBuilder或StringBuffer对象转换为字符串？
	\begin{options}
		\item toString()
		\item toStr()
		\item getString()
		\item getText()
	\end{options}
	\begin{answer}
		A
	\end{answer}
	\begin{explanation}
		\texttt{toString()} 方法返回字符串表示。
	\end{explanation}
\end{question}

\begin{question}
	以下关于Java中Math类的说法错误的是
	\begin{options}
		\item 提供了常用的数学计算方法
		\item 所有方法都是静态的
		\item 可以创建Math类的对象
		\item 包含了求绝对值、平方根等方法
	\end{options}
	\begin{answer}
		C
	\end{answer}
	\begin{explanation}
		\texttt{Math} 类是工具类，构造方法私有，不能实例化。
	\end{explanation}
\end{question}

\begin{question}
	以下哪个方法用于获取随机数？
	\begin{options}
		\item random()
		\item getRandom()
		\item Random.nextInt()
		\item Math.random()
	\end{options}
	\begin{answer}
		D
	\end{answer}
	\begin{explanation}
		\texttt{Math.random()} 返回 [0.0, 1.0) 的 double 随机数。
	\end{explanation}
\end{question}

\begin{question}
	以下关于Java中Calendar类的说法错误的是
	\begin{options}
		\item 用于操作日期和时间
		\item 是一个抽象类
		\item 可以直接创建对象
		\item 可以通过getInstance()方法获取实例
	\end{options}
	\begin{answer}
		C
	\end{answer}
	\begin{explanation}
		\texttt{Calendar} 是抽象类，需通过 \texttt{getInstance()} 获取实例。
	\end{explanation}
\end{question}

\begin{question}
	以下哪个方法用于格式化日期？
	\begin{options}
		\item format()
		\item parse()
		\item DateFormat.format
		\item SimpleDateFormat.format()
	\end{options}
	\begin{answer}
		D
	\end{answer}
	\begin{explanation}
		\texttt{SimpleDateFormat} 的 \texttt{format(Date)} 方法用于格式化。
	\end{explanation}
\end{question}

\begin{question}
	以下关于Java中Optional类的说法错误的是
	\begin{options}
		\item 用于避免空指针异常
		\item 可以包含值或为空
		\item 不支持值的获取和判断
		\item 提供了一些有用的方法
	\end{options}
	\begin{answer}
		C
	\end{answer}
	\begin{explanation}
		Optional 提供 \texttt{isPresent()}、\texttt{get()}、\texttt{orElse()} 等方法用于获取和判断值。
	\end{explanation}
\end{question}
